% SPDX-FileCopyrightText: 2026 Bart van de Lint
% SPDX-License-Identifier: AGPL-3.0-or-later
\section{Model detail against its effect}
\label{sec:model_detail}

\Cref{tab:model_detail} lists every part of the model with how it is modelled, how much it moves each rig's cost saving and which rig is cheapest, and its known error. A part the model can do without is switched off on the three ships of the development study (\cref{sec:res_inputs}): the calibration of the calm-water resistance, the added resistance in waves, the windage, the heel, leeway and rudder limits, and the rudder's drag. A part that development can change is moved over the ranges of its inputs in that study, one at a time; the propeller and the wind at height are not, so their effect is not measured here. \tbd{All at draft fidelity, where the routing noise of a kite's cost saving is up to 1.5 points (\cref{sec:checks}).}

\tbd{No part is detailed far beyond its effect. The heel, leeway and rudder limits move no cost saving by more than 0.4 points, and we keep them because they cost nothing and stop the router from using trims a ship could not hold. Two rounds of the size search lose at most 0.2 points against every size. The rudder drag costs rotors in a row up to 1.3 points. Towing the kite from the bow mast rather than the hull's lateral centre moves its cost saving by 0.1 points on average (\cref{sec:tow_point}). Each weather year alone moves a cost saving by up to 1.7 points, so a longer record would narrow the results without changing which rig is cheapest. Two parts change the cheapest rig or move a saving by several points and are coarser than their effect warrants, each a candidate for a better model: the extra wind a kite's launch needs in a sea, which moves its cost saving by up to 3.6 points and rests on one reported pair of wind speeds (a launch envelope measured at sea), and the corridor spacing (a finer or adaptive corridor), which moves the kite's cost saving by up to 1.8 points and decides between the kite and the rotors on the \SI{140}{m} ship, where they lie within 0.7 points of each other. \stale{router_grid}{On the corridor of the results the kite's saving still lies 1.3 points and the rotors' 0.5 below a reference finer in both, so the router is not converged and its error favours the rotors (\cref{sec:routing}).} The one mean power-law wind profile at the kite's height is coarse as well, and its effect is not measured here (better: hourly profiles from the ERA5 model levels).}

{\footnotesize
\setlength{\LTcapwidth}{\textwidth}
\setlength{\tabcolsep}{3.5pt}
\setlength{\extrarowheight}{1pt}
\begin{longtable}{@{}>{\raggedright\arraybackslash}p{2.5cm}>{\raggedright\arraybackslash}p{3.9cm}rrr>{\raggedright\arraybackslash}p{2.1cm}>{\raggedright\arraybackslash}p{3.5cm}@{}}
\caption{Each part of the model, how its effect is measured, its model, its effect on the cost saving of each rig, the ships where it changes which rig is cheapest, and its known error. The effect is the largest difference in a rig's cost saving, in percentage points, between the model as it stands and the model with the part changed, or with the inputs named at the part moved over their ranges, over the three ships of the development study. The parts switched off and the weather years come from the method study \studyrun{model_parts}, the router from \studyrun{convergence} and the size search from \studyrun{size_search}. \stale{convergence}{}}
\label{tab:model_detail}\\
\toprule
Part & Model & \multicolumn{3}{c}{Saving moved [pt]} & Cheapest rig changes & Known error \\
\cmidrule(lr){3-5}
 & & kite & wings & rotors & & \\
\midrule
\endfirsthead
\multicolumn{7}{@{}l}{\cref{tab:model_detail} (continued)}\\
\toprule
Part & Model & kite & wings & rotors & Cheapest rig changes & Known error \\
\midrule
\endhead
\bottomrule
\endfoot
\InputIfFileExists{sections/fidelity_table}{}{\multicolumn{7}{@{}l}{\tbd{Pending: run the fidelity\_table step.}}\\}
\end{longtable}}
